%!TEX root = ../main.tex

\begin{abstract}[it]

Negli ultimi anni la crescente diffusione di reti di sensori wireless (WSN) e di dispositivi IoT ha reso critico il problema della loro alimentazione: le soluzioni a batteria comportano infatti costi di sostituzione e manutenzione non trascurabili, specialmente quando i nodi sensore sono numerosi o di difficile accesso. L'\emph{energy harvesting}, ossia la conversione di energia ambientale -- vibrazioni meccaniche, calore, luce, segnali RF -- in energia elettrica utilizzabile, si propone come soluzione a questo problema, permettendo la realizzazione di nodi sensore autonomi e a manutenzione pressoché nulla.

Questo lavoro di tesi si apre con una panoramica delle principali tecnologie di trasduzione per l'harvesting da vibrazioni meccaniche -- piezoelettrica, elettromagnetica, elettrostatica e triboelettrica -- evidenziandone vantaggi e limiti, per poi concentrarsi sui raccoglitori a riluttanza variabile (\emph{Variable Reluctance Energy Harvester}, VREH). Questa tecnologia è particolarmente adatta alle applicazioni rotazionali a bassa velocità, poiché magnete e bobina restano entrambi fissi e la variazione di flusso concatenato è generata unicamente dal passaggio di una ruota dentata ferromagnetica. Viene quindi ripercorso lo stato dell'arte della ricerca su questa tecnologia, dai primi prototipi per il monitoraggio ferroviario fino alle topologie a polo ``a m'' ottimizzate per la densità di potenza volumetrica, sino al raccoglitore multifase (MP-VREH) a sei unità di pickup preso come riferimento in questo lavoro, pensato per l'integrazione in un cuscinetto di ruota di veicoli commerciali di grandi dimensioni, quali autobus, a velocità di rotazione comprese tra 100 e 400 giri/min (circa 20--80 km/h).

A partire dalla legge di Faraday viene poi sviluppato un modello elettrico completo del VREH, con la derivazione delle relazioni che legano la potenza estratta ai parametri geometrici ed elettrici del trasduttore -- numero di spire, resistenza di bobina, differenza di flusso tra posizione allineata e disallineata, numero di denti della ruota -- fino alla determinazione della potenza massima teoricamente trasferibile al carico e alla discussione delle condizioni di adattamento di impedenza in funzione della velocità di rotazione.

La parte centrale della tesi è dedicata alla progettazione dei circuiti di condizionamento della potenza per il raccoglitore a sei fasi considerato. Vengono proposte e confrontate diverse topologie di raddrizzamento, valutate dapprima tramite simulazione circuitale e successivamente tramite la realizzazione di un banco di prova sperimentale, comprensivo della progettazione di una scheda elettronica dedicata e della taratura di un sensore di corrente per la misura accurata della potenza estratta dal sistema. Infine, viene affrontata la scelta del circuito integrato di power management (PMIC) più adatto a regolare e immagazzinare l'energia raccolta, completando così la catena di conversione dall'energia meccanica rotazionale fino all'alimentazione stabile del carico elettronico finale.

\end{abstract}
